\section*{Overview}

Two pointers is one of the simplest ways to save a factor of \(n\). The pattern is to maintain a window or a pair of
positions and move each pointer only forward. If the state can be updated incrementally, the whole scan becomes linear.

\section*{When to Use It}

Reach for this pattern when:

\begin{itemize}
  \item the answer depends on a contiguous subarray or substring,
  \item the left and right boundaries only need to move forward,
  \item adding one element and removing one element can be maintained cheaply,
  \item a brute-force nested loop would revisit almost the same interval many times.
\end{itemize}

Classic tasks include longest valid subarray, shortest valid window, counting windows with some bound, or merging two
sorted lists.

\section*{Core Idea}

Maintain a current window \([l, r]\) and an invariant that describes when the window is valid. Expand \(r\) to add new
information; shrink \(l\) only when the invariant is violated or when you want to make the window minimal.

Because each pointer moves at most \(n\) times, the total number of operations is linear apart from the work needed to
maintain the window state.

\section*{Key Insight}

The technique works only when the window reacts monotonically enough to one-way movement.

For nonnegative arrays, if a sum is too large, moving the left pointer right can only decrease it. That makes the
state manageable. If negative numbers are allowed, the same trick often breaks and you need prefix sums or a different
data structure.

\section*{Operations / Main Technique}

\paragraph{Fixed-state sliding window.}
Add \(\texttt{a[r]}\), update counts or sums, then shrink while invalid.

\paragraph{Two sorted arrays.}
Use two indices to merge, count inversions in specialized settings, or match greedily.

\paragraph{Worked example.}
For the longest subarray with sum at most \(S\) and nonnegative values, expand the right end until the sum becomes too
large, then move the left end until the window is valid again. Every index enters once and leaves once.

\section*{Correctness Intuition}

At every step, the maintained invariant tells me what the current window means: valid, minimal valid, or maximal valid
ending at \(r\). Since the pointers never move backward, the algorithm never repeats a discarded state.

The one-way movement is the whole proof skeleton: each decision throws away configurations that will never become useful
again.

\section*{Complexity Analysis}

If both pointers only move forward, the total number of pointer updates is \(O(n)\). With \(O(1)\) amortized window
maintenance, the entire method is \(O(n)\).

\section*{Implementation}

The code includes:

\begin{itemize}
  \item a longest-window routine for nonnegative sums,
  \item a counter for subarrays with at most \(k\) distinct values.
\end{itemize}

Those two patterns cover a large fraction of contest sliding-window problems.

\section*{Common Pitfalls}

\begin{itemize}
  \item Applying the nonnegative-sum template when negative values are present.
  \item Forgetting whether the window is inclusive or half-open.
  \item Not updating the answer in the right place: before shrinking, after shrinking, or both.
  \item Letting the frequency map keep zero-count keys and then miscounting distinct values.
\end{itemize}

\section*{Variants / Extensions}

\begin{itemize}
  \item Count exactly \(k\) distinct by \(\texttt{at\_most(k) - at\_most(k - 1)}\).
  \item Monotone deque when the window also needs min or max queries.
  \item Two pointers on sorted arrays for pair sum, interval stabbing, or greedy matching.
  \item Meet-in-the-middle or prefix tricks when the window state is not one-way monotone.
\end{itemize}

\section*{Practice Problems}

\begin{itemize}
  \item Longest substring without repeated characters.
  \item Number of subarrays with at most \(k\) distinct values.
  \item Shortest subarray with sum at least \(S\) in the nonnegative setting.
  \item Merge-like matching between two sorted sequences.
\end{itemize}

\section*{References}

\begin{itemize}
  \item \href{https://usaco.guide/silver/two-pointers}{USACO Guide: Two Pointers}.
  \item \href{https://codeforces.com/catalog}{Codeforces Catalog}.
\end{itemize}
